\section{Rubrik Penilaian Komprehensif}

Penilaian komprehensif memerlukan pendekatan sistematis dan objektif. Rubrik yang baik mencakup \cite{stallings2017}:
\begin{itemize}
  \item \textbf{Kriteria penilaian}: Kriteria harus jelas, terukur, dan dapat diukur secara konsisten.
  \item \textbf{Feedback konstruktif}: Memberikan umpan balik yang mendalam dan konstruktif untuk perbaikan.
  \item \textbf{Authenticity}: Memastikan keaslian penilaian dan mencegah plagiarisme.
  \item \textbf{Fairness}: Menjamin proses penilaian yang adil dan transparan.
\end{itemize}


\subsection{Metode Penilaian}

Berbagai metode penilaian dapat digunakan untuk mengukur pencapaian:
\begin{itemize}
  \item \textbf{Observasi}: Pengamatan langsung proses pembelajaran dan partisipasi mahasiswa.
  \item \textbf{Portofolio}: Evaluasi karya mahasiswa untuk menunjukkan kemajuan praktikum.
  \item \textbf{Presentasi}: Presentasi hasil proyek atau penelitian.
  \item \textbf{Tes Tertulis}: Uji pemahaman dengan soal esai atau praktikum.
  \item \textbf{Peer Review}: Evaluasi kerja mahasiswa oleh teman sekelas.
  \item \textbf{Self-Assessment}: Refleksi diri dari pembelajaran.
\end{itemize}

Sesuai silabus: Tugas (15\%), Kuis (10\%), Praktikum pengenalan (20\%), Makalah mini (15\%), UTS (20\%), UAS (20\%).

\textbf{Kriteria nilai}: A (85--100), B (70--84), C (60--69), D (50--59), E (<50).

Pemetaan ringkas: Tugas/Kuis $\rightarrow$ Sub-CPMK 2.1--2.3, 1.1; Praktikum $\rightarrow$ 1.1--1.2; Makalah $\rightarrow$ CPMK-1 \& CPMK-2; UTS $\rightarrow$ CPMK-2 + 1.1--1.2; UAS $\rightarrow$ CPMK-1 \& CPMK-2 (termasuk 1.3--1.4).

